% Lösungen Einheit 2 von 3 – Bedingungen mit Ableitung (zusammenbau v0.1)
\einheitenkopf{Einheit 2 von 3 · Bedingungen mit Ableitung}

\erg{11}{a) zwei Gleichungen – $f(3) = 1$ und $f'(3) = 0$. \quad b) $c = -1$; $f(1) = a + b - 1 = -3$ und $f'(1) = 2a + b = 0$; $a = 2$, $b = -4$, also $f(x) = 2x^2 - 4x - 1$ \quad c) $h'(t) = 3at^2 + 2bt$; $h(4) = 64a + 16b = 64$ und $h'(4) = 48a + 8b = 24$; $a = -0{,}5$, $b = 6$, also $h(t) = -0{,}5t^3 + 6t^2$ \quad d) $c = 0$; $f(1) = 3 \cdot 1 + 1 = 4$ und $f'(1) = 3$: $a + b = 4$ und $2a + b = 3$; $a = -1$, $b = 5$, also $f(x) = -x^2 + 5x$ \quad e) Knickfrei in $S(0|2{,}5)$: $f(0) = r(0) = 2{,}5$ und $f'(0) = r'(0) = 0$, also $c = 2{,}5$, $b = 0$; $f(4) - g(4) = 1{,}2$ mit $g(4) = 0{,}5$: $16a + 2{,}5 = 1{,}7$, also $a = -0{,}05$ und $f(x) = -0{,}05x^2 + 2{,}5$ \quad f) Ansatz $p(x) = ax^2 + c$; $p(2) = 4a + c = 0$ und $p'(2) = 4a = \mathrm{tan}(-45^\circ) = -1$; $a = -0{,}25$, $c = 1$, also $p(x) = -0{,}25x^2 + 1$ und die Höhe $1$ m \quad g) Berührstelle: $6x - 3x^2 = 3$ liefert $x_0 = 1$; $f(x) = 3x^2 - x^3 + C$ mit $f(1) = t(1) = 5$: $2 + C = 5$, also $C = 3$ und $f(x) = -x^3 + 3x^2 + 3$ \quad h) $V'(t) = -20k \cdot e^{k \cdot t}$; $V'(0) = -20k = 4$, also $k = -0{,}2$ \quad i) Ansatz $p(x) = ax^4 + bx^2 + c$ mit $c = 1{,}8$; $p(2) = 16a + 4b + 1{,}8 = 0$ und $p'(2) = 32a + 4b = \mathrm{tan}(-45^\circ) = -1$; $a = 0{,}05$, $b = -0{,}65$, also $p(x) = 0{,}05x^4 - 0{,}65x^2 + 1{,}8$}
\erg{12}{a) $w(5) = 20 + b \cdot e^{5c} = 80$, also $b \cdot e^{5c} = 60$; $w'(5) = b \cdot c \cdot e^{5c} = 60c = -6$, also $c = -0{,}1$ und $b = 60 \cdot e^{0{,}5} \approx 98{,}92$}
\erg{13}{a) Der Hochpunkt liefert zwei Gleichungen; die Bedingung $f'(2) = 4a + b = 0$ fehlt, statt $b$ zu wählen. Richtig: $4a + 2b = 4$ und $4a + b = 0$, also $a = -1$, $b = 4$ und $f(x) = -x^2 + 4x - 2$. \quad b) $H$ liegt auf dem Graphen, das gibt $f(3) = 5$; im Hochpunkt ist die Tangente waagerecht, das gibt $f'(3) = 0$ – zwei verschiedene Aussagen, also zwei Gleichungen.}
